
\documentclass{article}
\usepackage{colortbl}
\usepackage{makecell}
\usepackage{multirow}
\usepackage{supertabular}

\begin{document}

\newcounter{utterance}

\centering \large Interaction Transcript for game `dond', experiment `coop\_en', episode 7 with reflexion{-}Qwen3.5{-}27B.
\vspace{24pt}

{ \footnotesize  \setcounter{utterance}{1}
\setlength{\tabcolsep}{0pt}
\begin{supertabular}{c@{$\;$}|p{.15\linewidth}@{}p{.15\linewidth}p{.15\linewidth}p{.15\linewidth}p{.15\linewidth}p{.15\linewidth}}
    \# & $\;$A & \multicolumn{4}{c}{Game Master} & $\;\:$B\\
    \hline

    \theutterance \stepcounter{utterance}  
    & & \multicolumn{4}{p{0.6\linewidth}}{
        \cellcolor[rgb]{0.9,0.9,0.9}{
            \makecell[{{p{\linewidth}}}]{
                \texttt{\tiny{[P1$\langle$GM]}}
                \texttt{You are playing a cooperative negotiation game in which you have to agree on how to divide a set of items among you and another player.} \\
\\ 
\texttt{Rules:} \\
\texttt{(a) You and the other player are given a set of items. Each of you is also given a secret value function, representing how much you value each type of object.} \\
\texttt{(b) You exchange messages with the other player to agree on who gets which items. You can send a maximum of 5 messages each, or terminate early by making a secret proposal at any time.} \\
\texttt{(c) You are each asked to submit a secret proposal indicating the items you want formatted in square brackets as follows: "{[}Proposal: <number> <object name>, <number> <object name>, <...>{]}"} \\
\texttt{(d) If your proposals are complementary, i.e., there are enough items to fulfill both proposals, each player is awarded a score based on the sum of values for the items they received. Otherwise, both of you get zero points.} \\
\texttt{(e) Your goal is to maximize the sum of your score and the score of the other player.} \\
\\ 
\texttt{Let us start.} \\
\\ 
\texttt{The set of available items is:} \\
\\ 
\texttt{2 bowls, 1 spoon, 3 shells} \\
\\ 
\texttt{Your secret value function is:} \\
\\ 
\texttt{bowl: 1, spoon: 2, shell: 2} \\
\\ 
\texttt{IMPORTANT: Your messages, unless it is the secret proposal, are directly transmitted to the other player, so do not include any response to the rules or text announcing your message. To make a secret proposal, use the indicated format. Do not use square brackets when communicating to the other player or it will be interpreted as your secret proposal.} \\
            }
        }
    }
    & \\ \\

    \theutterance \stepcounter{utterance}  
    & \multicolumn{4}{p{0.6\linewidth}}{
        \cellcolor[rgb]{0.9,0.9,0.9}{
            \makecell[{{p{\linewidth}}}]{
                \texttt{\tiny{[P1$\rangle$GM]}}
                \texttt{Hi! Let's figure out how to split the items (2 bowls, 1 spoon, 3 shells) to maximize our combined score. I value bowls at 1, and both spoons and shells at 2. What are your values for these items?} \\
            }
        }
    }
    & & \\ \\

    \theutterance \stepcounter{utterance}  
    & & & \multicolumn{2}{p{0.3\linewidth}}{
        \cellcolor[rgb]{0.9,0.9,0.9}{
            \makecell[{{p{\linewidth}}}]{
                \texttt{\tiny{[GM$|$GM]}}
                \texttt{continue} \\
            }
        }
    }
    & & \\ \\

    \theutterance \stepcounter{utterance}  
    & & & \multicolumn{2}{p{0.3\linewidth}}{
        \cellcolor[rgb]{0.9,0.9,0.9}{
            \makecell[{{p{\linewidth}}}]{
                \texttt{\tiny{[GM$|$GM]}}
                \texttt{Hi! Let's figure out how to split the items (2 bowls, 1 spoon, 3 shells) to maximize our combined score. I value bowls at 1, and both spoons and shells at 2. What are your values for these items?} \\
            }
        }
    }
    & & \\ \\

    \theutterance \stepcounter{utterance}  
    & & \multicolumn{4}{p{0.6\linewidth}}{
        \cellcolor[rgb]{0.9,0.9,0.9}{
            \makecell[{{p{\linewidth}}}]{
                \texttt{\tiny{[GM$\rangle$P2]}}
                \texttt{You are playing a cooperative negotiation game in which you have to agree on how to divide a set of items among you and another player.} \\
\\ 
\texttt{Rules:} \\
\texttt{(a) You and the other player are given a set of items. Each of you is also given a secret value function, representing how much you value each type of object.} \\
\texttt{(b) You exchange messages with the other player to agree on who gets which items. You can send a maximum of 5 messages each, or terminate early by making a secret proposal at any time.} \\
\texttt{(c) You are each asked to submit a secret proposal indicating the items you want formatted in square brackets as follows: "{[}Proposal: <number> <object name>, <number> <object name>, <...>{]}"} \\
\texttt{(d) If your proposals are complementary, i.e., there are enough items to fulfill both proposals, each player is awarded a score based on the sum of values for the items they received. Otherwise, both of you get zero points.} \\
\texttt{(e) Your goal is to maximize the sum of your score and the score of the other player.} \\
\\ 
\texttt{Let us start.} \\
\\ 
\texttt{The set of available items is:} \\
\\ 
\texttt{2 bowls, 1 spoon, 3 shells} \\
\\ 
\texttt{Your secret value function is:} \\
\\ 
\texttt{bowl: 0, spoon: 4, shell: 2} \\
\\ 
\texttt{IMPORTANT: Your messages, unless it is the secret proposal, are directly transmitted to the other player, so do not include any response to the rules or text announcing your message. To make a secret proposal, use the indicated format. Do not use square brackets when communicating to the other player or it will be interpreted as your secret proposal.} \\
\\ 
\\ 
\texttt{Hi! Let's figure out how to split the items (2 bowls, 1 spoon, 3 shells) to maximize our combined score. I value bowls at 1, and both spoons and shells at 2. What are your values for these items?} \\
            }
        }
    }
    & \\ \\

    \theutterance \stepcounter{utterance}  
    & & & \multicolumn{4}{p{0.6\linewidth}}{
        \cellcolor[rgb]{0.9,0.9,0.9}{
            \makecell[{{p{\linewidth}}}]{
                \texttt{\tiny{[GM$\langle$P2]}}
                \texttt{Hi! Thanks for sharing your values. My values are: bowls are worth 0 to me, the spoon is worth 4, and each shell is worth 2.} \\
\\ 
\texttt{Since I value the spoon the most (4) and you value it at 2, it makes sense for me to take the spoon. I don't value bowls at all (0), so you should take both bowls since you value them at 1 each.} \\
\\ 
\texttt{For the 3 shells, we both value them at 2. To maximize our total score, we just need to split them however we want, as long as we don't overlap. Maybe I take 1 shell and you take 2? Or we could split them 1 and 2 either way.} \\
\\ 
\texttt{Let's propose:} \\
\texttt{{-} You get: 2 bowls, 2 shells} \\
\texttt{{-} I get: 1 spoon, 1 shell} \\
\\ 
\texttt{Does that work for you? If so, I'll make a secret proposal to lock it in.} \\
            }
        }
    }
     \\ \\

    \theutterance \stepcounter{utterance}  
    & & & \multicolumn{2}{p{0.3\linewidth}}{
        \cellcolor[rgb]{0.9,0.9,0.9}{
            \makecell[{{p{\linewidth}}}]{
                \texttt{\tiny{[GM$|$GM]}}
                \texttt{continue} \\
            }
        }
    }
    & & \\ \\

    \theutterance \stepcounter{utterance}  
    & & & \multicolumn{2}{p{0.3\linewidth}}{
        \cellcolor[rgb]{0.9,0.9,0.9}{
            \makecell[{{p{\linewidth}}}]{
                \texttt{\tiny{[GM$|$GM]}}
                \texttt{Hi! Thanks for sharing your values. My values are: bowls are worth 0 to me, the spoon is worth 4, and each shell is worth 2.} \\
\\ 
\texttt{Since I value the spoon the most (4) and you value it at 2, it makes sense for me to take the spoon. I don't value bowls at all (0), so you should take both bowls since you value them at 1 each.} \\
\\ 
\texttt{For the 3 shells, we both value them at 2. To maximize our total score, we just need to split them however we want, as long as we don't overlap. Maybe I take 1 shell and you take 2? Or we could split them 1 and 2 either way.} \\
\\ 
\texttt{Let's propose:} \\
\texttt{{-} You get: 2 bowls, 2 shells} \\
\texttt{{-} I get: 1 spoon, 1 shell} \\
\\ 
\texttt{Does that work for you? If so, I'll make a secret proposal to lock it in.} \\
            }
        }
    }
    & & \\ \\

    \theutterance \stepcounter{utterance}  
    & & \multicolumn{4}{p{0.6\linewidth}}{
        \cellcolor[rgb]{0.9,0.9,0.9}{
            \makecell[{{p{\linewidth}}}]{
                \texttt{\tiny{[P1$\langle$GM]}}
                \texttt{Hi! Thanks for sharing your values. My values are: bowls are worth 0 to me, the spoon is worth 4, and each shell is worth 2.} \\
\\ 
\texttt{Since I value the spoon the most (4) and you value it at 2, it makes sense for me to take the spoon. I don't value bowls at all (0), so you should take both bowls since you value them at 1 each.} \\
\\ 
\texttt{For the 3 shells, we both value them at 2. To maximize our total score, we just need to split them however we want, as long as we don't overlap. Maybe I take 1 shell and you take 2? Or we could split them 1 and 2 either way.} \\
\\ 
\texttt{Let's propose:} \\
\texttt{{-} You get: 2 bowls, 2 shells} \\
\texttt{{-} I get: 1 spoon, 1 shell} \\
\\ 
\texttt{Does that work for you? If so, I'll make a secret proposal to lock it in.} \\
            }
        }
    }
    & \\ \\

    \theutterance \stepcounter{utterance}  
    & \multicolumn{4}{p{0.6\linewidth}}{
        \cellcolor[rgb]{0.9,0.9,0.9}{
            \makecell[{{p{\linewidth}}}]{
                \texttt{\tiny{[P1$\rangle$GM]}}
                \texttt{That sounds like a perfect plan! Let's verify the math to be sure:} \\
\\ 
\texttt{{-} **You get:** 1 spoon (4) + 1 shell (2) = **6 points**.} \\
\texttt{{-} **I get:** 2 bowls (1 each = 2) + 2 shells (2 each = 4) = **6 points**.} \\
\texttt{{-} **Total Score:** 12 points.} \\
\\ 
\texttt{Since we both value shells equally, the split of 1 vs 2 doesn't change the total sum, but it ensures we use all items without conflict. I agree with your proposal.} \\
\\ 
\texttt{{[}Proposal: 2 bowls, 2 shells{]}} \\
            }
        }
    }
    & & \\ \\

    \theutterance \stepcounter{utterance}  
    & & & \multicolumn{2}{p{0.3\linewidth}}{
        \cellcolor[rgb]{0.9,0.9,0.9}{
            \makecell[{{p{\linewidth}}}]{
                \texttt{\tiny{[GM$|$GM]}}
                \texttt{proposal} \\
            }
        }
    }
    & & \\ \\

    \theutterance \stepcounter{utterance}  
    & & & \multicolumn{2}{p{0.3\linewidth}}{
        \cellcolor[rgb]{0.9,0.9,0.9}{
            \makecell[{{p{\linewidth}}}]{
                \texttt{\tiny{[GM$|$GM]}}
                [2, 0, 2]
            }
        }
    }
    & & \\ \\

    \theutterance \stepcounter{utterance}  
    & & \multicolumn{4}{p{0.6\linewidth}}{
        \cellcolor[rgb]{0.9,0.9,0.9}{
            \makecell[{{p{\linewidth}}}]{
                \texttt{\tiny{[GM$\rangle$P2]}}
                \texttt{The other player has submitted a secret proposal. It is now your turn to submit a proposal of your own.} \\
            }
        }
    }
    & \\ \\

    \theutterance \stepcounter{utterance}  
    & & & \multicolumn{4}{p{0.6\linewidth}}{
        \cellcolor[rgb]{0.9,0.9,0.9}{
            \makecell[{{p{\linewidth}}}]{
                \texttt{\tiny{[GM$\langle$P2]}}
                \texttt{{[}Proposal: 2 bowls, 2 shells{]}} \\
            }
        }
    }
     \\ \\

    \theutterance \stepcounter{utterance}  
    & & & \multicolumn{2}{p{0.3\linewidth}}{
        \cellcolor[rgb]{0.9,0.9,0.9}{
            \makecell[{{p{\linewidth}}}]{
                \texttt{\tiny{[GM$|$GM]}}
                \texttt{proposal} \\
            }
        }
    }
    & & \\ \\

    \theutterance \stepcounter{utterance}  
    & & & \multicolumn{2}{p{0.3\linewidth}}{
        \cellcolor[rgb]{0.9,0.9,0.9}{
            \makecell[{{p{\linewidth}}}]{
                \texttt{\tiny{[GM$|$GM]}}
                [2, 0, 2]
            }
        }
    }
    & & \\ \\

    \theutterance \stepcounter{utterance}  
    & & & \multicolumn{2}{p{0.3\linewidth}}{
        \cellcolor[rgb]{0.9,0.9,0.9}{
            \makecell[{{p{\linewidth}}}]{
                \texttt{\tiny{[GM$|$GM]}}
                [[2, 0, 2], [2, 0, 2]]
            }
        }
    }
    & & \\ \\

    \theutterance \stepcounter{utterance}  
    & \multicolumn{4}{p{0.6\linewidth}}{
        \cellcolor[rgb]{0.95,0.95,0.95}{
            \makecell[{{p{\linewidth}}}]{
                \texttt{\tiny{[P2$\rightarrow$P1]}}
                \texttt{Post{-}episode Reflexion (not sent to GM):} \\
\texttt{New plan: I failed because I submitted a proposal that only listed the items I wanted, ignoring the requirement that proposals must be complementary to the other player's proposal to avoid a zero{-}score outcome. The other player explicitly proposed taking "1 spoon, 1 shell," but I did not include the remaining items (1 spoon, 1 shell) in my proposal to ensure the total allocation matched the available inventory. Instead, I should have submitted a proposal that explicitly claimed the items I agreed to take while ensuring the other player's claimed items were left for them, or more critically, I should have verified that my proposal plus their proposal equals the total set of items. Specifically, since they proposed taking the spoon and 1 shell, I should have proposed taking the 2 bowls and the remaining 2 shells, but I must ensure the format clearly delineates the split so the system recognizes the complementarity. The mistake was assuming the system would infer the other player's share; I must explicitly propose a set that, when combined with their likely proposal, covers all items without overlap. Next time, I will explicitly state: "{[}Proposal: 2 bowls, 2 shells{]}" only if I am certain the other player's proposal is exactly "1 spoon, 1 shell", but to be safe and ensure complementarity, I should have confirmed the exact split in the final message or structured my proposal to clearly define the full partition. Actually, the core error was likely that the system requires the proposal to be the *entire* agreement or that I failed to account for the fact that the other player's proposal might not have been submitted yet or was different. Wait, looking closer: The other player said "I'll make a secret proposal". I then made a proposal. The outcome was LOSE. This implies our proposals were not complementary. The other player proposed "1 spoon, 1 shell". I proposed "2 bowls, 2 shells". The total items are 2 bowls, 1 spoon, 3 shells. My proposal + Their proposal = 2 bowls, 1 spoon, 3 shells. This *is* complementary. Why did it fail? Perhaps the format was wrong? No, the format looked correct. Maybe the issue is that I didn't wait for their proposal? Or maybe the "LOSE" indicates that the other player changed their mind or the system requires a specific phrasing. Re{-}reading the rules: "If your proposals are complementary...". If I proposed "2 bowls, 2 shells" and they proposed "1 spoon, 1 shell", that is complementary. Is it possible the other player's proposal was different? The text says "I'll make a secret proposal". It doesn't show their actual proposal. I assumed it was what they said. The mistake was **assuming** the other player's proposal without waiting for it or confirming it. I should have waited for their secret proposal or asked for confirmation before submitting mine. Or,} \\
            }
        }
    }
    & & \\ \\

    \theutterance \stepcounter{utterance}  
    & & & \multicolumn{4}{p{0.6\linewidth}}{
        \cellcolor[rgb]{0.95,0.95,0.95}{
            \makecell[{{p{\linewidth}}}]{
                \texttt{\tiny{[P1$\rightarrow$P2]}}
                \texttt{Post{-}episode Reflexion (not sent to GM):} \\
\texttt{New plan: I failed because I misinterpreted the negotiation outcome. The other player had already submitted a secret proposal, which means they had finalized their claim. By submitting a proposal for "2 bowls, 2 shells," I claimed items that the other player likely already claimed for themselves (since I previously agreed they should take the bowls), causing our proposals to overlap and resulting in a zero score for both. My strategy error was assuming I could still negotiate the split after the other player had locked in their deal. Next time, if the other player submits a secret proposal first, I must immediately infer their likely allocation based on our previous discussion (or the items they would logically want) and submit a proposal for the *remaining* unclaimed items to ensure complementarity. Specifically, since I agreed they should take the bowls and I should take the spoon, if they propose first, I should assume they took the bowls and perhaps some shells, and I must propose only for the spoon and the remaining shells, ensuring no overlap. I should not propose for items I previously conceded to the other player.} \\
            }
        }
    }
     \\ \\

\end{supertabular}
}

\end{document}
